%% LyX 2.4.4 created this file.  For more info, see https://www.lyx.org/.
%% Do not edit unless you really know what you are doing.
\documentclass[english]{article}
\usepackage[T1]{fontenc}
\usepackage[utf8]{inputenc}
\setcounter{secnumdepth}{-2}
\setlength{\parindent}{0bp}
\usepackage{amsmath}
\usepackage{amsthm}
\usepackage{amssymb}

\makeatletter
%%%%%%%%%%%%%%%%%%%%%%%%%%%%%% Textclass specific LaTeX commands.
\theoremstyle{definition}
\newtheorem*{problem*}{\protect\problemname}

\@ifundefined{date}{}{\date{}}
%%%%%%%%%%%%%%%%%%%%%%%%%%%%%% User specified LaTeX commands.
\usepackage{graphicx}
\usepackage{geometry}
\newcommand{\limi}{\underset{n\rightarrow\infty}{\lim}}
\newcommand{\limxz}{\underset{x\rightarrow x_{0}}{\lim}}
\newcommand{\limfxz}{\underset{x\rightarrow x_{0}}{\lim}f(x)}
\newcommand{\limfxm}{\underset{x\rightarrow x_{0}^{-}}{\lim}f(x)}
\newcommand{\limfxp}{\underset{x\rightarrow x_{0}^{+}}{\lim}f(x)}
\newcommand{\limxm}{\underset{x\rightarrow x_{0}^{-}}{\lim}}
\newcommand{\limxp}{\underset{x\rightarrow x_{0}^{+}}{\lim}}
\newcommand{\sqrm}{M_{m\times m}(\mathbb{F})}
\newcommand{\seqa}{(a_{n})_{n=1}^{\infty}}
\newcommand{\seqx}{(x_{n})_{n=1}^{\infty}}
\newcommand{\funcdr}{f:D\rightarrow\mathbb{R}}
\newcommand{\der}{\limxz\frac{f(x)-f(x_{0})}{x-x_{0}}}
\usepackage[all]{pdfprivacy}

\makeatother

\usepackage{babel}
\providecommand{\problemname}{Problem}
\begin{document}
\title{{\Large Semester 2026A, Exam B, Q1\vspace{-2em}}}

\maketitle
\begin{problem*}
Let $A,B\in M_{n\times n}(\mathbb{F})$ be such
that $AB=I_{n}$. Prove that $A$ and $B$ are inverses of each other.
\end{problem*}

\medskip{}

\rule[0.5ex]{1\columnwidth}{1pt}

\medskip{}

\begin{proof}
As we know, the homogeneous system $Bx=0$ has at least one solution
(the trivial solution), hence we can mark the set of its solutions
as $U$, and definitionally $\ker(T_{B})=U$. 

Let $x\in U$, and let us multiply both sides on the left by A:

\[
A\cdot Bx=A\cdot0\therefore\quad A\cdot Bx=0
\]

matrix multiplication is associative, hence 
\[
AB\cdot x=0\therefore\quad I_{n}\cdot x=0\therefore\quad x=0
\]

Which means $\ker(T_{B})=\{0\}$.

Recall the theorem that $T_{B}$ is injective if and only if its kernel
is 0, which means $T_{B}:\mathbb{\mathbb{F}}^{n}\rightarrow\mathbb{\mathbb{F}}^{n}$
is injective.

Furthermore, recall that a linear transformation between finite dimensional
vector spaces of the same dimension is surjective if and only if it
is injective. $\mathbb{\mathbb{F}}^{n}$ is a finite dimensional vector
space of dimension $n\in\mathbb{N}$, therefore $T_{B}$ is bijective,
which means $B$ is invertible, therefore B has an inverse $B^{-1}$
and $BB^{-1}=B^{-1}B=I_{n}$.

If we prove that $B^{-1}=A$, we get $BA=AB=I_{n}$, as required. 

\[
I_{n}B^{-1}=B^{-1}\underset{\text{substitution}}{\therefore}(AB)B^{-1}=B^{-1}\underset{\text{associativity}}{\therefore}A\left(BB^{-1}\right)=B^{-1}
\]

\[
\therefore AI_{n}=B^{-1}\therefore A=B^{-1}
\]
\end{proof}
\begin{quote}
This is an exam problem I translated and solved. The original exam
was written by the Linear Algebra 1 course lecturers at HUJI.
\end{quote}

\end{document}
